% ============================================================
%  ACCT 270 — Group Assignment 2: Predicting Accounting Fraud
% ============================================================
\documentclass[12pt]{article}

\usepackage[margin=1in]{geometry}
\usepackage{array}
\usepackage{booktabs}
\usepackage{amsmath}
\usepackage{amssymb}
\usepackage{multirow}
\usepackage{caption}
\usepackage{setspace}
\usepackage{hyperref}
\usepackage{xcolor}
\usepackage{colortbl}
\usepackage{listings}
\usepackage{lmodern}
\usepackage{microtype}
\usepackage{parskip}
\usepackage{fancyhdr}
\usepackage{threeparttable}

% ---- formatting ----
\setstretch{1.15}
\definecolor{siggreen}{RGB}{198,239,206}  % light green for significant cells

\pagestyle{fancy}
\fancyhf{}
\rhead{ACCT 270 $\cdot$ Group Assignment 2}
\lhead{Predicting Accounting Fraud}
\cfoot{\thepage}

% ---- SAS code style ----
\lstdefinestyle{sas}{
  basicstyle=\ttfamily\footnotesize,
  commentstyle=\color{gray}\itshape,
  breaklines=true,
  frame=single,
  numbers=left,
  numberstyle=\tiny\color{gray},
  numbersep=5pt,
  showstringspaces=false,
  tabsize=4,
  captionpos=b
}

\hypersetup{colorlinks=true, linkcolor=black, citecolor=black, urlcolor=blue}

% helper: green cell for significant coefficients
\newcommand{\sig}[1]{\cellcolor{siggreen}#1}

% ============================================================
\begin{document}

% ---- Header ----
\begin{center}
  {\large\bfseries Group 2}\\[0.3em]
  {\large\bfseries Predicting Accounting Fraud}\\[1em]
  \begin{tabular}{rl}
    \textbf{Team Name:}       & XXX \\
    \textbf{Individual Names:}& XXX \\
  \end{tabular}
\end{center}

\bigskip

% ============================================================
% EQUATION
% ============================================================

We estimate the following fraud prediction model:

\begin{equation}
  \text{Fraud}_{t} = \alpha
    + \theta_1\,\text{M-Flag}
    + \theta_2\,\text{F-Flag}
    + \beta_1 Z_1
    + \beta_2 Z_2
    + \cdots
    + \beta_n Z_n
    + \varepsilon
  \label{eq:fraud}
\end{equation}

\noindent
$\text{Fraud}_t$ is an indicator variable equal to one if the financial results for
the firm-year were the subject of an SEC Accounting and Auditing Enforcement Release
(\textsc{aaer}), and zero otherwise.
M-Flag is the flag for fraud using the Beneish M-score (M-score $> -1.78$).
F-Flag is the flag for fraud using the Dechow F-score (F-score $> 1.0$).
$Z$'s are additional predictor variables described below.
Table~1 provides variable definitions.

% ============================================================
% TABLE 1: VARIABLE DEFINITIONS
% ============================================================

\begin{table}[h!]
\centering
\caption*{\textbf{Table 1: Variable Definitions}}
\begin{tabular}{p{2.2cm} p{12.5cm}}
\toprule
\textbf{Variable} & \textbf{Definition} \\
\midrule
\multicolumn{2}{l}{\textsc{dependent variable}} \\[3pt]
\textbf{Fraud$_t$} &
  Indicator $= 1$ if the firm-year is subject to an SEC Accounting and Auditing
  Enforcement Release (\textsc{aaer}); 0 otherwise. Source: Dechow et al.\ (2011)
  \textsc{aaer} database, matched to Compustat via \texttt{gvkey} and fiscal year.\\[8pt]

\multicolumn{2}{l}{\textsc{independent variables}} \\[3pt]
\textbf{M-FLAG} &
  $= 1$ if M-Score $> -1.78$; 0 otherwise.\\
  & M-Score $= -4.84 + 0.92\,\text{DSR} + 0.528\,\text{GMI} + 0.404\,\text{AQI}
    + 0.892\,\text{SGI} + 0.115\,\text{DEPI} - 0.172\,\text{SGAI}
    + 4.679\,\text{ACCRUALS} - 0.327\,\text{LEVI}$ \\[2pt]
  & $\text{DSR} = \dfrac{RECT_t/SALE_t}{RECT_{t-1}/SALE_{t-1}}$\quad
    $\text{GMI} = \dfrac{1-(COGS_{t-1}/SALE_{t-1})}{1-(COGS_t/SALE_t)}$\\[6pt]
  & $\text{AQI} = \dfrac{1-((PPENT_t+ACT_t)/AT_t)}{1-((PPENT_{t-1}+ACT_{t-1})/AT_{t-1})}$\quad
    $\text{SGI} = \dfrac{SALE_t}{SALE_{t-1}}$\\[6pt]
  & $\text{DEPI} = \dfrac{(DP_{t-1}-AM_{t-1})/(DP_{t-1}-AM_{t-1}+PPENT_{t-1})}{(DP_t-AM_t)/(DP_t-AM_t+PPENT_t)}$\\[6pt]
  & $\text{SGAI} = \dfrac{XSGA_t/SALE_t}{XSGA_{t-1}/SALE_{t-1}}$\quad
    $\text{ACCRUALS} = \dfrac{IB_t - OANCF_t}{AT_t}$\quad
    $\text{LEVI} = \dfrac{(LCT_t+DLTT_t)/AT_t}{(LCT_{t-1}+DLTT_{t-1})/AT_{t-1}}$\\[4pt]
  & \textit{Missing DLTT, XSGA, AM set to zero. If any index denominator is zero,
    that index is set to 1.}\\[8pt]

\textbf{F-FLAG} &
  $= 1$ if F-Score $> 1.0$; 0 otherwise.\\
  & $\hat{y} = -7.893 + 0.790\,\text{RSST} + 2.518\,\text{CH\_REC}
    + 1.191\,\text{CH\_INV} + 1.979\,\text{SOFT}
    + 0.171\,\text{CH\_CS} - 0.932\,\text{CH\_ROA} + 1.029\,\text{ISSUE}$\\[2pt]
  & $\text{Prob} = e^{\hat{y}}/(1+e^{\hat{y}})$;\quad
    $\text{F-Score} = \text{Prob}/0.0037$\\[2pt]
  & $\text{RSST} = (IB_t - OANCF_t)/AT_t$\quad
    $\text{CH\_REC} = (RECT_t - RECT_{t-1})/AT_t$\\[3pt]
  & $\text{CH\_INV} = (INVT_t - INVT_{t-1})/AT_t$\quad
    $\text{SOFT} = (AT_t - PPENT_t - CHE_t)/AT_t$\\[3pt]
  & $\text{CH\_CS} = \dfrac{SALE_t-(RECT_t-RECT_{t-1})}{SALE_{t-1}-(RECT_{t-1}-RECT_{t-2})} - 1$\quad
    $\text{CH\_ROA} = IB_t/AT_t - IB_{t-1}/AT_{t-1}$\\[3pt]
  & $\text{ISSUE} = 1$ if $DLTIS + SSTK > 0$; 0 otherwise.
    \textit{Missing DLTIS, SSTK set to zero. If CH\_CS denominator is zero, CH\_CS set to 1.}\\[8pt]

\textbf{ACCRUAL} &
  $(IB_t - OANCF_t)/AT_t$. Total accruals scaled by total assets.
  High accruals signal a gap between reported earnings and cash flows — a classic
  indicator of earnings management. \textit{Predicted sign: positive.}\\[6pt]

\textbf{ASSET\_GROWTH} &
  $(AT_t - AT_{t-1})/AT_{t-1}$. Year-over-year growth in total assets.
  Aggressive expansion may reflect inflated capitalizations or acquisition-driven
  manipulation to sustain headline growth. \textit{Predicted sign: positive.}\\[6pt]

\textbf{LEV} &
  $(LCT_t + DLTT_t)/AT_t$. Total debt scaled by assets.
  Highly leveraged firms face covenant pressure and may misreport to avoid
  technical default or maintain credit access. \textit{Predicted sign: positive.}\\

\bottomrule
\end{tabular}
\end{table}

% ============================================================
% TRANSITION
% ============================================================

\noindent
Following Beneish (1999), Beneish, Lee, and Nichols (2011), and Dechow et al.\ (2011),
we examine the M-Flag and F-Flag variables as required predictors of fraud.
We supplement these with three additional variables — ACCRUAL, ASSET\_GROWTH, and LEV
— each motivated by a distinct fraud incentive channel.
Note: ROA $(= IB/AT)$ was initially included but dropped due to near-perfect collinearity
with ACCRUAL ($r = 0.905$, as both contain $IB/AT$). SALES\_GROWTH was dropped due to
extreme outliers (std~$= 61$, max~$= 11{,}880$) that rendered it uninformative.
Table~2 summarizes the descriptive statistics.

% ============================================================
% TABLE 2: DESCRIPTIVE STATISTICS
% ============================================================

\begin{table}[h!]
\centering
\caption*{\textbf{Table 2: Descriptive Statistics}}
\begin{threeparttable}
\begin{tabular}{lccccccc}
\toprule
\textbf{Variable} & \textbf{Obs} & \textbf{Mean} & \textbf{Std} &
  \textbf{25th} & \textbf{Median} & \textbf{75th} & \textbf{Years} \\
\midrule
Fraud$_t$    & 41,959 & 0.007 & 0.085 & 0.000 & 0.000 & 0.000 & 2003--2014 \\
M-Flag       & 41,959 & 0.121 & 0.326 & 0.000 & 0.000 & 0.000 & 2003--2014 \\
F-Flag       & 41,959 & 0.348 & 0.476 & 0.000 & 0.000 & 1.000 & 2003--2014 \\
ACCRUAL      & 41,959 & $-$0.079 & 0.508 & $-$0.099 & $-$0.053 & $-$0.019 & 2003--2014 \\
ASSET\_GROWTH & 41,959 & 0.126 & 0.646 & $-$0.039 & 0.051 & 0.166 & 2003--2014 \\
LEV          & 41,959 & 0.417 & 0.405 & 0.240 & 0.388 & 0.539 & 2003--2014 \\
\bottomrule
\end{tabular}
\begin{tablenotes}
\small\item
All variables defined in Table~1. Sample: fiscal years 2003--2014, requiring
positive $AT$, $SALE$, $PPENT$ and non-missing values for all model variables.
\end{tablenotes}
\end{threeparttable}
\end{table}

% ============================================================
% TABLE 3: ALL-YEARS REGRESSION
% ============================================================

Table~3 provides regression results. $t$-statistics are reported in parentheses.
*, **, *** denote coefficients significantly different from zero at the 10\%, 5\%,
and 1\% level respectively. Cells with significant coefficients are shaded green.

\begin{table}[h!]
\centering
\caption*{\textbf{Table 3: Regression Results — All Years}}
\begin{threeparttable}
\begin{tabular}{lc}
\toprule
\textbf{Independent Variable} & \textbf{ALL YEARS} \\
                               & \textbf{Fraud}$_t$ \\
\midrule
Intercept      & \sig{$0.0053^{***}$} \\
               & \sig{$(7.22)$}       \\[4pt]
M-Flag         & $0.0003$             \\
               & $(0.22)$             \\[4pt]
F-Flag         & \sig{$0.0073^{***}$} \\
               & \sig{$(8.31)$}       \\[4pt]
ACCRUAL        & $-0.0007$            \\
               & $(-0.59)$            \\[4pt]
ASSET\_GROWTH  & $0.0011$             \\
               & $(1.61)$             \\[4pt]
LEV            & $-0.0021$            \\
               & $(-1.45)$            \\[6pt]
Observations   & $41{,}959$           \\
\bottomrule
\end{tabular}
\begin{tablenotes}
\small\item
OLS estimates. Dependent variable: Fraud$_t$. $t$-statistics in parentheses.
$^{*}p<0.10$, $^{**}p<0.05$, $^{***}p<0.01$. Shaded cells are significant at 10\% or better.
\end{tablenotes}
\end{threeparttable}
\end{table}

\noindent
\textbf{Findings:}
Consistent with our predictions, we find that F-Flag is positive and highly significant
($\hat{\theta}_2 = 0.0073$, $t = 8.31$), indicating that firms flagged by the Dechow
F-Score are approximately 0.73 percentage points more likely to commit fraud — a
near-doubling of the unconditional fraud rate of 0.70\%.
M-Flag is positive but statistically insignificant ($t = 0.22$), consistent with
the finding in Dechow et al.\ (2011) that the F-Score has greater discriminatory
power than the M-Score for predicting \textsc{aaer} firms.
Among our additional variables, ACCRUAL and LEV are in the predicted direction but
insignificant, while ASSET\_GROWTH is marginally positive ($t = 1.61$, $p = 0.11$).
The low $R^2$ ($0.0018$) is expected given the rarity of fraud ($<1\%$ of firm-years)
and is consistent with values reported in the literature.

\newpage

% ============================================================
% TABLE 4: YEAR-BY-YEAR 2003--2008
% ============================================================

\begin{table}[h!]
\centering
\caption*{\textbf{Table 4: Regression Results by Year — 2003 to 2008}}
\begin{threeparttable}
\setlength{\tabcolsep}{5pt}
\begin{tabular}{lcccccc}
\toprule
\textbf{Independent} & \textbf{2003} & \textbf{2004} & \textbf{2005} &
  \textbf{2006} & \textbf{2007} & \textbf{2008} \\
\textbf{Variable}    & \textbf{Fraud}$_t$ & \textbf{Fraud}$_t$ & \textbf{Fraud}$_t$ &
  \textbf{Fraud}$_t$ & \textbf{Fraud}$_t$ & \textbf{Fraud}$_t$ \\
\midrule
Intercept
  & \sig{$0.0123^{***}$} & \sig{$0.0116^{***}$} & \sig{$0.0116^{***}$}
  & $0.0054^{*}$         & \sig{$0.0068^{**}$}  & \sig{$0.0066^{***}$} \\
  & \sig{$(3.43)$}       & \sig{$(3.30)$}        & \sig{$(3.27)$}
  & $(1.94)$             & \sig{$(2.52)$}        & \sig{$(3.05)$}        \\[4pt]
M-Flag
  & $-0.0059$    & $-0.0011$    & $0.0003$    & $0.0006$    & $-0.0008$   & $0.0005$ \\
  & $(-0.99)$    & $(-0.21)$    & $(0.05)$    & $(0.14)$    & $(-0.20)$   & $(0.13)$ \\[4pt]
F-Flag
  & \sig{$0.0210^{***}$} & \sig{$0.0109^{***}$} & \sig{$0.0101^{***}$}
  & \sig{$0.0089^{***}$} & $0.0048^{*}$         & $0.0013$               \\
  & \sig{$(4.91)$}       & \sig{$(2.98)$}        & \sig{$(2.82)$}
  & \sig{$(3.22)$}       & $(1.83)$              & $(0.54)$               \\[4pt]
ACCRUAL
  & $0.0024$    & $0.0037$    & $-0.0009$   & $0.0001$    & $0.0075$    & $0.0031$ \\
  & $(0.21)$    & $(0.50)$    & $(-0.07)$   & $(0.03)$    & $(0.91)$    & $(0.71)$ \\[4pt]
ASSET\_GROWTH
  & $0.0007$    & $0.0009$    & \sig{$0.0078^{**}$} & $0.0049$   & $0.0003$   & $0.0011$ \\
  & $(0.30)$    & $(0.29)$    & \sig{$(2.04)$}       & $(1.59)$   & $(0.15)$   & $(0.40)$ \\[4pt]
LEV
  & $-0.0053$   & $-0.0077$   & \sig{$-0.0139^{**}$} & $-0.0077$  & $-0.0058$  & $-0.0053$ \\
  & $(-0.78)$   & $(-1.16)$   & \sig{$(-1.99)$}       & $(-1.42)$  & $(-1.15)$  & $(-1.41)$ \\[6pt]
Observations
  & $4{,}152$   & $3{,}996$   & $3{,}803$   & $3{,}721$   & $3{,}529$   & $3{,}448$ \\
\bottomrule
\end{tabular}
\begin{tablenotes}
\small\item
OLS estimates by fiscal year. Dependent variable: Fraud$_t$.
$t$-statistics in parentheses. $^{*}p<0.10$, $^{**}p<0.05$, $^{***}p<0.01$.
Shaded cells are significant at 10\% or better.
\end{tablenotes}
\end{threeparttable}
\end{table}

% ============================================================
% TABLE 5: YEAR-BY-YEAR 2009--2014
% ============================================================

\begin{table}[h!]
\centering
\caption*{\textbf{Table 5: Regression Results by Year — 2009 to 2014}}
\begin{threeparttable}
\setlength{\tabcolsep}{5pt}
\begin{tabular}{lcccccc}
\toprule
\textbf{Independent} & \textbf{2009} & \textbf{2010} & \textbf{2011} &
  \textbf{2012} & \textbf{2013} & \textbf{2014} \\
\textbf{Variable}    & \textbf{Fraud}$_t$ & \textbf{Fraud}$_t$ & \textbf{Fraud}$_t$ &
  \textbf{Fraud}$_t$ & \textbf{Fraud}$_t$ & \textbf{Fraud}$_t$ \\
\midrule
Intercept
  & $0.0046^{*}$  & $0.0045^{*}$  & $0.0015$      & $-0.0002$     & $0.0010$      & $-0.0010$ \\
  & $(1.83)$      & $(1.95)$      & $(0.59)$       & $(-0.06)$     & $(0.53)$      & $(-0.65)$ \\[4pt]
M-Flag
  & $0.0005$    & $0.0011$    & $-0.0056$   & $-0.0011$   & $0.0005$    & $0.0037$ \\
  & $(0.12)$    & $(0.26)$    & $(-1.36)$   & $(-0.24)$   & $(0.14)$    & $(1.46)$ \\[4pt]
F-Flag
  & $0.0025$    & $0.0018$    & \sig{$0.0055^{**}$}  & \sig{$0.0097^{***}$} & $0.0024$   & \sig{$0.0044^{***}$} \\
  & $(0.83)$    & $(0.63)$    & \sig{$(2.13)$}        & \sig{$(3.39)$}        & $(1.08)$   & \sig{$(2.88)$}        \\[4pt]
ACCRUAL
  & $-0.0011$   & $0.0010$    & $0.0036$    & $-0.0002$   & $-0.0009$   & $-0.0032$ \\
  & $(-0.15)$   & $(0.32)$    & $(0.40)$    & $(-0.03)$   & $(-0.15)$   & $(-0.61)$ \\[4pt]
ASSET\_GROWTH
  & $0.0003$    & $0.0008$    & $-0.0001$   & $-0.0031$   & $-0.0006$   & \sig{$0.0035^{**}$} \\
  & $(0.25)$    & $(0.48)$    & $(-0.03)$   & $(-1.28)$   & $(-0.37)$   & \sig{$(2.25)$}       \\[4pt]
LEV
  & $0.0005$    & $0.0013$    & $0.0049$    & $0.0071$    & $0.0032$    & $0.0001$ \\
  & $(0.09)$    & $(0.29)$    & $(0.97)$    & $(1.27)$    & $(0.90)$    & $(0.03)$ \\[6pt]
Observations
  & $3{,}449$   & $3{,}297$   & $3{,}178$   & $3{,}166$   & $3{,}131$   & $3{,}089$ \\
\bottomrule
\end{tabular}
\begin{tablenotes}
\small\item
OLS estimates by fiscal year. Dependent variable: Fraud$_t$.
$t$-statistics in parentheses. $^{*}p<0.10$, $^{**}p<0.05$, $^{***}p<0.01$.
Shaded cells are significant at 10\% or better.
\end{tablenotes}
\end{threeparttable}
\end{table}

\noindent
\textbf{Findings:}
Results from estimating the regression by year indicate that F-Flag is the most
consistent predictor across the sample period, achieving significance at the 5\%
level or better in 8 of 12 years (2003--2006, 2011, 2012, 2014) and at the 10\%
level in 2007. Its coefficient is positive in every year, ranging from 0.0013
in 2008 to 0.0210 in 2003. M-Flag is never significant in any individual year
and alternates in sign, confirming its limited incremental predictive power
conditional on F-Flag. The 2008--2010 period shows no significant predictors,
consistent with reduced \textsc{aaer} enforcement activity during and after the
financial crisis, as evidenced by the decline in fraud rates from 1.57\% in 2003
to 0.44\% in 2008. ASSET\_GROWTH achieves significance in 2005 and 2014,
consistent with the prediction that aggressive asset expansion is episodically
associated with fraudulent reporting during periods of elevated capital market
activity.

% ============================================================
% APPENDIX: SAS CODE
% ============================================================

\newpage
\noindent\textbf{Appendix A}

\lstset{style=sas}
\begin{lstlisting}
/*=================================================================
  ACCT 270 - Group Assignment 2: Predicting Accounting Fraud
  Model: Fraud = a + b1*MFLAG + b2*FFLAG + b3*ACCRUAL
                   + b4*ASSET_GROWTH + b5*LEV + e
  Sample: 2003-2014 | Data: Compustat, AAER
=================================================================*/

libname mydata "s:\sasdata";
libname home   "d:\users\mcbai\desktop\programs";
options ps=max ls=max nocenter;

/*-----------------------------------------------------------------
  STEP 1: CONSTRUCT M-SCORE (Beneish 1999)
  MSCORE = -4.84 + 0.92*DSR + 0.528*GMI + 0.404*AQI + 0.892*SGI
           + 0.115*DEPI - 0.172*SGAI + 4.679*ACCRUALS - 0.327*LEVI
  MFLAG = 1 if MSCORE > -1.78
-----------------------------------------------------------------*/

/* Examine AAER dataset to confirm structure and year coverage */
proc contents data=mydata.aaer1982to2014; run;
proc freq data=mydata.aaer1982to2014; tables fyear; run;

/* Pull Compustat; require positive AT, SALE, PPENT and
   non-missing core variables for M-score and F-score           */
data variables;
    set mydata.compustat1990to2019;
    if AT > 0 and SALE > 0 and PPENT > 0;
    if nmiss(cogs, dp, rect, act, lct, ib, oancf) = 0;

    /* Fill optional variables with zero if missing */
    if dltt = . then dltt = 0;
    if xsga = . then xsga = 0;
    if am   = . then am   = 0;
    if invt = . then invt = 0;
    if che  = . then che  = 0;
    if dltis= . then dltis= 0;
    if sstk = . then sstk = 0;

    /* M-Score raw components */
    DS      = RECT / SALE;
    GM      = 1 - (COGS / SALE);
    AQ      = 1 - ((PPENT + ACT) / AT);
    SG      = SALE;
    DEP     = (DP - AM) / (DP - AM + PPENT);
    SGA     = XSGA / SALE;
    ACCRUAL = (IB - OANCF) / AT;          /* Also used as Z1 predictor */
    LEV     = (LCT + DLTT) / AT;          /* Also used as Z3 predictor */

    /* ROA computed but excluded from regression:
       near-perfect collinearity with ACCRUAL (r=0.905) */
    ROA = IB / AT;
run;

/* Self-join to prior year for M-score index ratios,
   lagged asset/sale values for growth variables          */
proc sql;
    create table addlagyear as
    select a.*,
           b.DS as DS_m1, b.GM as GM_m1, b.AQ as AQ_m1,
           b.SG as SG_m1, b.DEP as DEP_m1, b.SGA as SGA_m1,
           b.ACCRUAL as ACCRUAL_m1, b.LEV as LEV_m1,
           b.AT as AT_m1, b.IB as IB_m1,
           b.RECT as RECT_m1, b.INVT as INVT_m1,
           b.SALE as SALE_m1
    from variables a
    inner join variables b
        on a.gvkey = b.gvkey and a.fyear - 1 = b.fyear;
quit;

/* Compute M-Score indices (set to 1 if denominator = 0),
   MSCORE, MFLAG, and additional predictor variables       */
data Mscore;
    set addlagyear;
    if DS_m1  > 0 then DSR  = DS  / DS_m1;  else DSR  = 1;
    if GM     > 0 then GMI  = GM_m1 / GM;   else GMI  = 1;
    if AQ_m1  > 0 then AQI  = AQ  / AQ_m1; else AQI  = 1;
    if SG_m1  > 0 then SGI  = SG  / SG_m1; else SGI  = 1;
    if DEP    > 0 then DEPI = DEP_m1 / DEP; else DEPI = 1;
    if SGA_m1 > 0 then SGAI = SGA / SGA_m1; else SGAI = 1;
    if LEV_m1 > 0 then LEVI = LEV / LEV_m1; else LEVI = 1;

    MSCORE = -4.84 + 0.92*DSR + 0.528*GMI + 0.404*AQI + 0.892*SGI
             + 0.115*DEPI - 0.172*SGAI + 4.679*ACCRUAL - 0.327*LEVI;
    if MSCORE > -1.78 then MFLAG = 1; else MFLAG = 0;

    /* Z2: asset growth using lagged total assets */
    if AT_m1   > 0 then ASSET_GROWTH = (AT   - AT_m1)   / AT_m1;
    /* SALES_GROWTH dropped: extreme outliers (std=61, max=11880) */
    if SALE_m1 > 0 then SALES_GROWTH = (SALE - SALE_m1) / SALE_m1;
run;

/*-----------------------------------------------------------------
  STEP 2: CONSTRUCT F-SCORE (Dechow et al. 2011)
  Requires 2-year lagged receivables for CH_CS
  FFLAG = 1 if FSCORE > 1.0
-----------------------------------------------------------------*/

/* Get 2-year lagged receivables from raw Compustat */
proc sql;
    create table addlagyear2 as
    select a.*, b.rect as rect_m2
    from Mscore a
    inner join mydata.compustat1990to2019 b
        on a.gvkey = b.gvkey and a.fyear - 2 = b.fyear;
quit;

/* Compute 7 F-Score components, log-odds, probability, F-Score */
data Fscore;
    set addlagyear2;
    if nmiss(rect,rect_m1,rect_m2,ib,ib_m1,invt,invt_m1,
             che,at,at_m1) = 0;

    RSST   = ACCRUAL;
    CH_REC = (RECT - RECT_m1) / AT;
    CH_INV = (INVT - INVT_m1) / AT;
    SOFT   = (AT - PPENT - CHE) / AT;
    CH_ROA = (IB/AT) - (IB_m1/AT_m1);

    /* If denominator of CH_CS is zero, set to 1 */
    if SALE_m1 - (RECT_m1 - RECT_m2) > 0
        then CH_CS = ((SALE-(RECT-RECT_m1)) /
                      (SALE_m1-(RECT_m1-RECT_m2))) - 1;
        else CH_CS = 1;

    if DLTIS + SSTK > 0 then ISSUE = 1; else ISSUE = 0;

    PredictedValue = -7.893 + 0.790*RSST + 2.518*CH_REC
                     + 1.191*CH_INV + 1.979*SOFT
                     + 0.171*CH_CS - 0.932*CH_ROA + 1.029*ISSUE;
    /* Cap at 709 to prevent exp() overflow */
    if PredictedValue > 709 then PredictedValue = 709;

    Probability = exp(PredictedValue) / (1 + exp(PredictedValue));
    Fscore      = Probability / 0.0037;
    if FSCORE > 1 then FFLAG = 1; else FFLAG = 0;
run;

/*-----------------------------------------------------------------
  STEP 3: MERGE AAER, CODE FRAUD, RESTRICT TO 2003-2014
-----------------------------------------------------------------*/

/* Left join: keep all firm-years; match to AAER => fraud=1 */
proc sql;
    create table addAAER as
    select a.*, b.P_aaer
    from FSCORE a
    left join mydata.aaer1982to2014 b
        on a.gvkey = b.gvkey and a.fyear = b.fyear;
quit;

data Fraud;
    set addAAER;
    if 2003 <= fyear <= 2014;
    if nmiss(mflag, fflag) = 0;
    if nmiss(accrual, asset_growth, lev) = 0;
    if not missing(P_aaer) then FRAUD = 1; else FRAUD = 0;
run;

/*-----------------------------------------------------------------
  STEP 4: DESCRIPTIVE STATISTICS
-----------------------------------------------------------------*/

proc means data=fraud n mean std p25 p50 p75 maxdec=3;
    var fraud mflag fflag accrual asset_growth lev;
quit;

proc freq data=fraud; tables fyear; run;

proc means data=fraud mean maxdec=4;
    var fraud; class fyear;
run;

proc corr data=fraud;
    var mflag fflag accrual asset_growth lev;
quit;

/*-----------------------------------------------------------------
  STEP 5: PART I - POOLED REGRESSION (2003-2014)
-----------------------------------------------------------------*/

proc reg data=fraud;
    model fraud = mflag fflag accrual asset_growth lev;
    title 'Part I: Pooled Fraud Regression (2003-2014)';
quit;

/*-----------------------------------------------------------------
  STEP 6: PART II - YEAR-BY-YEAR REGRESSIONS
  BY statement loops over sorted fyear values
-----------------------------------------------------------------*/

proc sort data=fraud; by fyear; run;

proc reg data=fraud;
    by fyear;
    model fraud = mflag fflag accrual asset_growth lev;
    title 'Part II: Year-by-Year Fraud Regressions';
quit;
\end{lstlisting}

\end{document}
